\subsection{QUOTIENT SPACES AND PRODUCT SPACES OF NORMED SPACES}

PROPOSITION 9. Let E be a normed space over a non-discrete valued division ring K, and let H be a closed vector subspace of E. If, for each class \(\dot{\mathbf{x}}\in\mathbf{E}/\mathbf{H}\) , we put \(||\dot{\mathbf{x}}||=\inf_{\mathbf{x}\in\dot{\mathbf{x}}}\|\mathbf{x}||\) , the function \(||\dot{\mathbf{x}}||\) is a norm on the vector space E/H,

and the topology defined by this norm is the quotient by H of the topology of E. By Remark 2 of no. 1, \(d(\dot{\boldsymbol{x}}, \dot{\boldsymbol{y}}) = ||\dot{\boldsymbol{x}} - \dot{\boldsymbol{y}}||\) is an invariant metric on E/H which defines the quotient by H of the topology of E. It remains only to show that \(||t\dot{\boldsymbol{x}}|| = |t|.||\dot{\boldsymbol{x}}||\) , and this follows immediately from the definition of \(||\dot{\boldsymbol{x}}||\) {[}Chapter IV, § 5, no. 7, formula (23){]}.

The norm \(||\dot{\pmb{x}}||\) may also be interpreted as follows: it is the distance (in E) of every point \(\pmb{x} \in \dot{\pmb{x}}\) from the subspace H, for the points of \(\dot{\pmb{x}}\) are the points \(\pmb{x} - \pmb{z}\) , where \(\pmb{z}\) runs through H.

PROPOSITION 10. Let \((\mathbf{E}_i)_{1 \leqslant i \leqslant n}\) be a finite family of normed spaces over a non-discrete valued division ring K, and let \(\mathbf{E} = \prod_{i=1}^{n} \mathbf{E}_i\) be the product vector space. If for each \(x = (x_i) \in \mathbf{E}\) we put \(||x|| = \sup_{1 \leqslant i \leqslant n} ||x_i||\) , then the function \(||x||\) is a norm on E, and the topology it defines on E is the product of the topologies of the \(\mathbf{E}_i\) .

For if \(\pmb{x} = (\pmb{x}_i)\) and \(\pmb{y} = (\pmb{y}_i)\) , we have \(\pmb{x} + \pmb{y} = (\pmb{x}_i + \pmb{y}_i)\) and therefore

\[
\begin{array}{r l} | | \boldsymbol {x} + \boldsymbol {y} | | = & \sup _ {i} | | \boldsymbol {x} _ {i} + \boldsymbol {y} _ {i} | | \leqslant \sup _ {i} (| | \boldsymbol {x} _ {i} | | + | | \boldsymbol {y} _ {i} | |) \\ & \leqslant \sup _ {i} | | \boldsymbol {x} _ {i} | | + \sup _ {i} | | \boldsymbol {y} _ {i} | | = | | \boldsymbol {x} | | + | | \boldsymbol {y} | |. \end{array}
\]

On the other hand, it is clear that \(||tx|| = |t|.||x||\) , and that if \(||x|| = 0\) , then \(||x_i|| = 0\) and therefore \(x_i = 0\) for \(1 \leqslant i \leqslant n\) , so that \(x = 0\) ; hence \(||x||\) is a norm on E. Also the relation \(||x|| < a\) is equivalent to the \(n\) relations \(||x_i|| < a\) , and therefore the norm \(||x||\) defines the product topology on E.

Similarly we can show that the functions \(\sum_{i=1}^{n} \|x_i\|\) and \(\sqrt{\sum_{i=1}^{n} \|x_i\|^2}\) are norms on E; the inequalities (2) show that all three norms are equivalent.

In particular, in the (left or right) vector space \(K^{n}\) , if we put

\[
p _ {1} (\boldsymbol {x}) = \sup _ {i} | x _ {i} |, \quad p _ {2} (\boldsymbol {x}) = \sum_ {i = 1} ^ {i = 1} | x _ {i} |, \quad p _ {3} (\boldsymbol {x}) = \sqrt {\sum_ {i = 1} ^ {n} | x _ {i} | ^ {2}}
\]

for \(\pmb{x} = (x_{i})_{1\leqslant i\leqslant n}\) , then the three functions \(p_1, p_2, p_3\) are equivalent norms which define on \(\mathbf{K}^n\) the topology which is the product of the topologies of the factors \(\mathbf{K}\) .

